\section{Introducción}

El presente trabajo articula una ontología del tiempo extramental discreto fundamentada en el marco del \textbf{Reloj Extramental Discreto} (RECD) y en la noción de \textbf{presente estratificado}. A diferencia de las concepciones newtonianas del tiempo como flujo continuo o de las aproximaciones puramente subjetivistas, el RECD propone que el tiempo extramental posee una estructura \textit{discreta y ordinal} detectable a través de conjunciones significativas en la configuración de persistencia de sistemas dinámicos complejos \citep{padilla_bayesian_accumulator_2026,padilla_sintesis_magna_2026}.

La investigación empírica previa ha validado el RECD en regímenes caóticos (mapas logísticos acoplados con $r=3.8$), en el atractor de Lorenz sometido a cambios de parámetro controlados y en series epidemiológicas reales. Estos estudios han puesto de manifiesto que la \textit{intensificación de la persistencia} local, medida mediante hiper-persistencia y trapping estructural, no se traduce automáticamente en reorganización global del sistema. La clave de la transición radica en la \textbf{antisincronización} entre componentes: la desviación estándar de la escala de persistencia local $\tau_s$.

El presente artículo desarrolla esta observación en una ontología coherente. Se sostiene que el presente es un fenómeno \textbf{estratificado y fragmentado}, organizado en tres capas jerárquicas cuya integración está regulada por la antisincronización. Esta última actúa como un mecanismo de resistencia que mantiene la autonomía de los presentes locales e impide que intensificaciones aisladas produzcan efectos estructurales a nivel del sistema.

La contribución principal consiste en:
\begin{itemize}
    \item Una definición rigurosa y operacional de las tres capas del presente.
    \item La caracterización de la antisincronización como regulador de la permeabilidad entre capas.
    \item La explicitación de una dinámica de resistencia bidireccional (de abajo hacia arriba y de arriba hacia abajo).
    \item La reformulación de la memoria como \textit{retención activa} sustentada en trapping e hiper-persistencia.
\end{itemize}

El artículo se estructura de la siguiente manera. En la Sección 2 se exponen los fundamentos técnicos del RECD y se definen con precisión los constructos centrales. La Sección 3 resume la evidencia empírica que sustenta la ontología propuesta. Las Secciones 4 y 5 desarrollan la estructura de capas y la dinámica de resistencia. La Sección 6 aborda la noción de memoria activa y sus implicaciones para el tiempo extramental. Finalmente, se discuten limitaciones, la analogía con sistemas de múltiples escalas y se señalan líneas abiertas de investigación.